% Pursuing Stacks (cotes 134-2 to 134-8): archive cote 134-6, site page 114 (PDF page 115), typescript page 432.
% typescript. Draft compared: ps5.tex:7452-7528.
% AI-transcribed by Claude Opus 5.5 (claude-opus-5-5, adaptive thinking, effort high), Message Batches API, batch ps3 (msgbatch_01XdFmuVVBAcgc7EFxdtCqsj), 2026-11-12; not yet checked by eye.
\scanpage{134-6/114}
\mspage{432}

surrounding the familiar homology and cohomology notations H. and H$^{\bullet}$.\hand{\unsure{.}}
This flavor is at least partially preserved, it seems to me, in the
notation $\mathrm{R}\Gamma_{A}$ (because of the geometric intuition tied with the sections
functor), whereas there is not yet a familiar geometric notion of a
``cosections functor''. As we would like to have the duality symmetry
reflected as perfectly as possible in the notation, I am going to
use the notations

\xout{(123)}
\[
(123)\quad \left\{ \begin{array}{l}
\mathrm{RH}^{\bullet}(A,F^{\bullet}) = \mathrm{R}\varprojlim_{A^{o}}\ (F^{\bullet}) \quad ( = \mathrm{R}\Gamma_{A}(F^{\bullet})\hand{)} \ : \quad D^{+}(A^{\wedge}_{\hand{\mathcal{M}}}) \hand{\to \mathcal{M}} \\
\mathrm{LH}_{\bullet}(A,F_{\bullet}) = \mathrm{L}\varinjlim_{A^{o}}\ (F_{\bullet}) \quad : \quad D^{-}(A^{\wedge}_{\hand{\mathcal{N}}}) \hand{\to \mathcal{N}} \qquad .
\end{array} \right.
\]

With these notations, the duality isomorphisms (120,120') take the
form (as announced) :
\[
(124)\quad \left\{ \begin{array}{l}
(\mathrm{RH}^{\bullet}(A,F^{\bullet}))^{o} \simeq \mathrm{LH}_{\bullet}(B,(F^{\bullet})^{o}) \\
(\mathrm{LH}_{\bullet}(A,F_{\bullet}))^{o} \simeq \mathrm{RH}^{\bullet}(B,(F_{\bullet})^{o}) \qquad ,
\end{array} \right.
\]
where the first members are defined in terms of cohomology resp\hand{.} homolo-%
-gy invariants with respect to the pair $(A,\hand{\mathcal{M}})$ resp. $(A,\hand{\mathcal{N}})$ , whereas
the second members denote homology resp. cohomology invariants with
respect to the dual pairs $(B,\hand{\mathcal{N}})$ resp. $(B,\hand{\mathcal{M}})$.

\uline{Remarks} This perfect symmetry, or rather essential identity, between
``homology'' or ``integration'' and ``cohomology'' or ``cointegration'', is
obtained here at the price of working with presheaves with values in
rather general abelian categories,, subjected\xout{mainly} to some simple exact-%
ness properties. It should be remembered moreover that for the time
being, RH$^{\bullet}$ has not been defined as the total right derived functor of
the sections or \xout{lim} inverse limits functor, therefore the notations
(121) and (121') are somewhat misleading. To feel really at ease, we
should \xout{really develop} still \inserted{work out} conditions which ensure that $A^{\wedge}_{\hand{\mathcal{M}}}$ say has
enough injectives and that RHom$_{k}$ can be defined also using such resolu-%
tions - in which case we'll expect too to have a satisfactory forma-%
lism for the \uline{RHom}$_{k}$ functor.

\margin{114}
